\section{Exact descent and the height-two map}\label{h3:sec:descent}

The coheight-one coefficient comparison first gives two actual maps at
height two: the unit and the determinant class.  We then use those maps
in the finiteness argument needed for the ordinary comparison.
Throughout these sections \(p\geq5\).  Put \(U=u^{-1}\), so that
\(|U|=2\) and \(v_1=xU^{p-1}\).

\subsection{Relative descent and exact functors}

The tests below are applied to already formed, completed Morava
cooperation spectra.  Their preservation of descent is a finite
categorical statement, which we establish first.  We then identify the
relative cooperation terms and their actual cofaces.
A commutative algebra is \emph{descendable} if the unit belongs to the
thick tensor ideal it generates.

\begin{lemma}[Exact Amitsur descent]\label{h3:lem:exact-amitsur}
Let $\mathcal C$ be a stable symmetric monoidal $\infty$-category with
limits and with tensor product exact in each variable.  Let $E$ be a
descendable commutative algebra, write $\mathbf1$ for the unit, and put
$I_E=\fib(\mathbf1\to E)$.  For its full Amitsur object
$C^q=E^{\otimes(q+1)}$, the partial totalization
\[
 \operatorname{Tot}_s C^\bullet
   =\lim_{[q]\in\Delta_{\leq s}}C^q
\]
is a finite cubical limit.  Its augmentation fiber, including its map
to the unit, is
\begin{equation}\label{h3:des:partial-amitsur-fiber}
 \bigl(\fib(\mathbf1\longrightarrow\operatorname{Tot}_s C^\bullet)
       \longrightarrow\mathbf1\bigr)
 \simeq \bigl(I_E^{\otimes(s+1)}\longrightarrow\mathbf1\bigr).
\end{equation}
The transition between successive fibers applies $I_E\to\mathbf1$
to the last factor.  There is an integer $N$ such that every composite
of $N$ successive maps between these augmentation fibers is null.
Every exact functor from
$\mathcal C$ to a stable category preserves the augmented
totalization $\mathbf1\simeq\operatorname{Tot}C^\bullet$ and this
same null-composite bound.
\end{lemma}

\begin{proof}
Here $\Delta_{\leq s}$ is the full simplex category on
$[0],\ldots,[s]$, so its limit includes the codegeneracies.
Let $\mathcal P_s^\times$ be the finite poset of nonempty subsets of
$\{0,\ldots,s\}$.  The functor
\[
 \theta_s:\mathcal P_s^\times\longrightarrow\Delta_{\leq s},
 \qquad J\longmapsto[|J|-1],
\]
sends an inclusion to the order-preserving injection between the
ordered subsets.  It is initial for limits.  To verify the cofinality
criterion, fix $[q]$, where $q\leq s$.  An object of
$\theta_s\downarrow[q]$ is a nonempty subset $J$ together with a
weakly increasing map $J\to\{0,\ldots,q\}$; its morphisms are
extensions of these partial maps.  Thus this comma category is the
nonempty face poset of the simplicial complex $P(s,q)$ whose simplices
are the lists
\[
 (i_0,j_0),\ldots,(i_b,j_b),\qquad
 0\leq i_0<\cdots<i_b\leq s,
 \quad 0\leq j_0\leq\cdots\leq j_b\leq q.
\]
Its nerve is the barycentric subdivision of $P(s,q)$.

The complex $P(s,q)$ is contractible whenever $s\geq q$.
The base case $P(0,0)$ is a point.  For the induction, the simplices
containing the vertex $(s,j)$, together with their faces, form the
cone on $P(s-1,j)$.
Start with the cone with apex $(s,q)$ on $P(s-1,q)$; this includes
all simplices with first coordinates less than $s$ and is
contractible even if its base is not.  Attach the cones with apices
$(s,j)$, for $j=q-1,\ldots,0$.  The intersection of each such cone
with the preceding union is exactly its base $P(s-1,j)$: a simplex
cannot contain two vertices with first coordinate $s$.
Every attaching base is contractible by induction, since
$j<q\leq s$ implies $j\leq s-1$.  These are inclusions of finite
simplicial subcomplexes, hence cofibrations; attaching a cone along a
contractible subcomplex preserves contractibility.  This proves the
claim, including the endpoint $q=s$, where no assertion about the
contractibility of the initial base $P(s-1,s)$ was used.  The comma
categories are therefore contractible, proving initiality of
$\theta_s$.

It follows that $\operatorname{Tot}_s C^\bullet$ is the limit of the
punctured finite $(s+1)$-cube
\[
 J\longmapsto\bigotimes_{i=0}^s
   \begin{cases}E,&i\in J,\\ \mathbf1,&i\notin J,\end{cases}
 \qquad \varnothing\ne J\subseteq\{0,\ldots,s\},
\]
whose maps insert units.  Successive fibers and exactness of tensor
product identify its augmentation fiber as the arrow in
\eqref{h3:des:partial-amitsur-fiber}.  Omitting the last cubical
direction gives the stated transition.  Thus the model and its
transition refer to full partial totalizations, not a coface-only
truncation.  An exact functor preserves this finite cubical limit.
The finite-cube formula appears in \cite[Proposition~2.14]{MNN2017};
the comparison with full partial totalizations and the augmentation
fiber have been proved here in the stated setting.

For a descendable $E$, the partial Amitsur tower represents the
constant pro-object with value the unit
\cite[Proposition~3.20]{Mathew}.  Its limit is preserved by every
finite-limit-preserving functor \cite[Proposition~3.10]{Mathew}.
Together with the finite cubical comparison, this proves the
asserted identification of the transformed totalization.  There is
also an integer $N$ such that every composite of $N$ maps which
become null after tensoring with $E$ is null
\cite[Proposition~3.27]{Mathew}.  One can read this bound from a finite
thick-ideal construction of the unit: it is one for $E$, is unchanged
by retracts and tensor products, and adds across a cofiber sequence.
The arrow $I_E\to\mathbf1$ becomes null after tensoring with $E$:
its tensor is the composite through the null map $I_E\to E$,
followed by multiplication.  Thus the error transitions in
\eqref{h3:des:partial-amitsur-fiber} become null after tensoring with
$E$, and their $N$-fold composites are null.  An exact functor
preserves those null composites, supplying the same convergence
bound after applying the functor.
\end{proof}

\begin{proposition}[Exact relative Morava descent]
\label{h3:prop:exact-descent}
Let \(n\geq1\), let \(k\) be a finite field containing \(F=\F_{p^n}\), and let
\(R_{n,k}=L_{K(n)}S_k\).  In the \(K(n)\)-local category of
\(S_k\)-modules, the algebra \(E_n(k)\) is descendable over the unit
\(R_{n,k}\).  Write \(\widehat\otimes\) for the tensor product in this
local category and put
\[
 \mathcal N_n^q=E_n(k)^{\widehat\otimes_{R_{n,k}}(q+1)}
 \qquad(q\geq0).
\]
Its augmented Amitsur object totalizes to \(R_{n,k}\).  Every exact
functor from this local category to a stable category preserves that
totalization and a uniform nilpotence bound for its error tower.
In particular this applies to the inclusion in spectra followed by
ordinary smash, a finite Moore quotient, a chromatic localization, or
a monochromatic fiber.

There are identifications of graded rings
\begin{equation}\label{h3:des:relative-cooperations}
 \pi_*\mathcal N_n^q
    \cong C_{\hthreects}(P_n^q,E_n(k)_*),
\end{equation}
with the maximal-ideal topology on the coefficients.  Under these
identifications the actual cofaces and codegeneracies induce those of
the continuous action cobar construction, and the augmented unit is
the constant unit.  Each \(\mathcal N_n^q\) is an ordinary
\(E_n(k)\)-module through its first factor.  This is the module
structure used when an ordinary exact functor is applied to a term.
\end{proposition}

\begin{proof}
We begin with the standard coefficient field, keeping the scalar
extension separate.  Set \(R_n=L_{K(n)}S\),
\(E_F=E_n(F)\), \(\Gamma_n=\Gal(F/\F_p)\), and
\(G_n^0=P_n\rtimes\Gamma_n\).  The Hopkins--Ravenel theorem in the
form of \cite[Theorem~4.18]{Mathew} says that \(E_F\) generates the unit
\(L_n\mathbb S_{(p)}\) as a thick tensor ideal in the \(E(n)\)-local
category.  Applying the symmetric monoidal localization \(L_{K(n)}\)
preserves this finite generation statement
\cite[Corollary~3.21]{Mathew}.  The map \(\mathbb S_{(p)}\to S\) is a
mod-\(p\) equivalence, hence a \(K(n)\)-equivalence for \(n\geq1\),
so the resulting unit is \(R_n\).  The same localized base case is
stated directly in \cite[Proposition~10.10]{Mathew} for Morava theory
attached to a perfect residue field; the relative assertion still needs
the following scalar step.  Now base change along
\(R_n\to R_{n,k}\) inside the \(K(n)\)-local category.  Thus
\[
 B=E_F\widehat\otimes_{R_n}R_{n,k}
\]
is descendable over \(R_{n,k}\).

We identify one factor of this base change and prove that it is
itself descendable.  A \(\Z_p\)-basis of \(W(k)\) gives an equivalence
\(R_{n,k}\simeq R_n^{\vee [k:\F_p]}\) of underlying \(R_n\)-modules.
More precisely, the natural map \(R_n\otimes_S S_k\to R_{n,k}\) is
an equivalence: its source is a finite sum of \(K(n)\)-local copies of
\(R_n\), and localizing the finite free decomposition of \(S_k\) gives
the same map.  The natural scalar map therefore identifies the graded coefficient
ring of \(B\) with \(E_{F,*}\otimes_{\Z_p}W(k)\); no infinite limit is
interchanged with scalar extension here.  The finite étale algebra
of constants splits as
\begin{equation}\label{h3:des:embedding-idempotents}
 W(F)\otimes_{\Z_p}W(k)
   \xrightarrow{\ \cong\ }
   \prod_{\sigma:F\hookrightarrow k}W(k),
 \qquad a\otimes b\longmapsto(\sigma(a)b)_\sigma.
\end{equation}
Let \(e_\sigma\) denote its embedding idempotents.  Idempotent
localization splits \(B\) into the finite product of
\(B_\sigma=B[e_\sigma^{-1}]\): on homotopy groups the map to this
product is the displayed orthogonal-idempotent decomposition, hence
is an equivalence.  Each factor is even periodic, with degree-zero
ring \(W(k)[[u_1,\ldots,u_{n-1}]]\).  Its orientation is the universal
Honda deformation after the indicated unramified base change.  We may
therefore take the factor for the chosen inclusion \(\iota:F\subseteq k\)
as the model \(E_n(k)\) used here.

For \(\gamma\in\Gamma_n\), the Galois automorphism of the standard
factor \(E_F\), base changed with the identity on \(R_{n,k}\), is an
\(R_{n,k}\)-algebra automorphism of \(B\).  It permutes the idempotents
by \(e_\sigma\mapsto e_{\sigma\gamma^{-1}}\).  In particular it
identifies all the factors \(B_\sigma\) as \(R_{n,k}\)-algebras while
the second scalar action stays fixed.  This is the factor equivalence
being used; coefficient Frobenius on the chosen \(E_n(k)\) alone is
semilinear over \(W(k)\).  The underlying module of \(B\) is consequently
a finite sum of modules equivalent to \(E_n(k)\).  The thick tensor
ideal generated by \(E_n(k)\) contains \(B\), and the ideal generated
by \(B\) contains the unit.  This proves descendability of the chosen
factor.

Apply \Cref{h3:lem:exact-amitsur} in the \(K(n)\)-local category of
\(S_k\)-modules, with its local tensor product and unit \(R_{n,k}\).
The inclusion into spectra is exact, as are the ordinary tests in the
statement.  The lemma therefore proves the asserted totalization and
uniform error-tower bound for the relative Amitsur object and for each
of those tests.  We now keep the same \(B\) and its embedding
idempotents to identify the terms and maps of that object.

It remains to prove \eqref{h3:des:relative-cooperations} with its
maps.  The completed-cooperations theorem of Devinatz--Hopkins applies
to the standard \(E_F\) and the standard extended group \(G_n^0\):
\cite[Proposition~2.2 and formulas~(2.5)--(2.7)]{DevinatzHopkins}
identifies
\[
 \pi_*\bigl(E_F^{\widehat\otimes_{R_n}(q+1)}\bigr)
     \cong C_{\hthreects}((G_n^0)^q,E_{F,*}).
\]
The tensor on the left is the \(K(n)\)-local tensor; with its local
unit \(R_n\) it is the completed smash appearing in that theorem.
Base change gives a canonical equivalence of actual terms
\[
 B^{\widehat\otimes_{R_{n,k}}(q+1)}
 \simeq
 E_F^{\widehat\otimes_{R_n}(q+1)}
       \widehat\otimes_{R_n}R_{n,k}.
\]
Because \(R_{n,k}\) is finite free over \(R_n\), the preceding
coefficient identification becomes
\begin{equation}\label{h3:des:basechanged-cooperations}
 \pi_*\bigl(B^{\widehat\otimes_{R_{n,k}}(q+1)}\bigr)
 \cong C_{\hthreects}((G_n^0)^q,E_{F,*})\otimes_{\Z_p}W(k)
 \cong C_{\hthreects}((G_n^0)^q,B_*).
\end{equation}
The last isomorphism uses a finite basis of \(W(k)\), so it only
commutes continuous functions with a finite direct sum.

We spell out the coordinate convention because it controls the first
coface and the field component.  Write the action of \(G_n^0\) on
\(B\) for its action on \(E_F\) tensored with the identity on
\(R_{n,k}\).  In inhomogeneous coordinates the evaluation of a class
in the left side of \eqref{h3:des:basechanged-cooperations} at
\((g_1,\ldots,g_q)\) is induced by the actual multiplication map
\begin{equation}\label{h3:des:inhomogeneous-evaluation}
 \mu\circ\bigl(1\widehat\otimes g_1\widehat\otimes
       (g_1g_2)\widehat\otimes\cdots\widehat\otimes
       (g_1\cdots g_q)\bigr):
 B^{\widehat\otimes_{R_{n,k}}(q+1)}\longrightarrow B.
\end{equation}
For \(q=0\) this is the identity.  For \(q\geq1\), the evaluation in
Devinatz--Hopkins applies
\(h_i^{-1}\) to the first \(q\) factors and leaves the last factor
fixed.  Put \(h_i=g_i\cdots g_q\) and then apply \(h_1\) to the value.
This continuous change of variables gives exactly
\eqref{h3:des:inhomogeneous-evaluation}.  Their natural evaluation
formulas apply to every homotopy class, so this argument does not assume
that products of factor coefficients generate the completed term.

The projection \(B\to E_n(k)=B_\iota\) is a unital algebra map.
In degree \(q\), the induced projection to \(\mathcal N_n^q\) is the
all-\(\iota\) idempotent summand of the tensor product.  Under
\eqref{h3:des:inhomogeneous-evaluation}, its projector acts on a function
by multiplication by
\[
 e_\iota\,g_1(e_\iota)\,(g_1g_2)(e_\iota)\cdots
                 (g_1\cdots g_q)(e_\iota).
\]
The action of \(\Gamma_n\) on the embedding idempotents is free and
transitive, and its kernel in \(G_n^0\) is \(P_n\).  This product is
\(e_\iota\) precisely when every \(g_i\) lies in \(P_n\), and is zero
otherwise.  Since \(P_n^q\) is clopen in \((G_n^0)^q\), the selected
summand is exactly \(C_{\hthreects}(P_n^q,(B_\iota)_*)\).  The
\(P_n\)-action on the selected coefficient ring fixes its \(W(k)\)
scalars.  This proves \eqref{h3:des:relative-cooperations}.

The same evaluation maps identify all cobar maps.  For a coefficient
function \(f\) in degree \(q\), their formulas are
\[
 \begin{aligned}
 (d^0f)(g_1,\ldots,g_{q+1})
     &=g_1\bigl(f(g_2,\ldots,g_{q+1})\bigr),\\
 (d^if)(g_1,\ldots,g_{q+1})
     &=f(g_1,\ldots,g_i g_{i+1},\ldots,g_{q+1})
       &&(1\leq i\leq q),\\
 (d^{q+1}f)(g_1,\ldots,g_{q+1})
     &=f(g_1,\ldots,g_q),\\
 (s^if)(g_1,\ldots,g_{q-1})
     &=f(g_1,\ldots,g_i,1,g_{i+1},\ldots,g_{q-1})
       &&(0\leq i<q).
 \end{aligned}
\]
They follow by inserting a unit or multiplying adjacent factors in
\eqref{h3:des:inhomogeneous-evaluation}.  The unital projection from
the \(B\)-Amitsur object to the \(E_n(k)\)-Amitsur object commutes with
these maps.  Thus its selected cofaces are obtained by projecting the
actual cofaces and then restricting the displayed formulas to \(P_n\).
Their higher compatibilities are those of that actual Amitsur object.
The augmented unit evaluates to the constant unit.  The first coface is
\(R_{n,k}\)-linear and has the displayed varying stabilizer action on
coefficients; the other cofaces and all codegeneracies are linear for
the first-factor \(E_n(k)\)-module structures.  This proves the claims
about maps and naturality.
\end{proof}

The notation \(C_{\hthreects}(P_n^q,E_n(k))\) below refers to the actual
term \(\mathcal N_n^q\), equipped with
\eqref{h3:des:relative-cooperations}; it does not infer an equivalence
of spectra from a coefficient isomorphism.  The inclusion of the
\(K(n)\)-local category in spectra is exact, but its tensor product is
still the completed one.  Only after forming a term do we regard it as
an ordinary module and apply an ordinary exact functor.

We also record the finite-field descent, including its unit.  For
\(\tau\in\Gal(k/\F_p)\), combine the standard automorphism
\(\tau|_F\) on \(E_F\) with \(\tau\) on the second factor
\(R_{n,k}\).  On constants, evaluation at \(\iota\) sends
\((\tau|_F)(a)\otimes\tau(b)\) to \(\tau(\iota(a)b)\).
Consequently this combined action preserves \(B_\iota\), is semilinear
over \(W(k)\), and conjugates \(P_n\) by the usual Honda action.
It gives the coefficient-field action on the relative object just
constructed and respects its augmentation and cofaces.

Let \(\Gamma=\Gal(k/\F_p)\).  A normal-basis element \(\bar a\in k\)
has nonzero trace: otherwise the sum of its conjugates would be a
nontrivial \(\F_p\)-linear dependence.  Lift it to \(a_0\in W(k)\).
The conjugates of \(a_0\) are a \(\Z_p\)-basis, since their determinant
is a unit modulo \(p\), and its trace is a unit.  Replacing \(a_0\) by
\(a=a_0/\operatorname{Tr}(a_0)\) gives a normal basis with trace one.
The map
\[
 \operatorname{Ind}_{1}^{\Gamma}S=\bigvee_{\gamma\in\Gamma}S
       \longrightarrow S_k,
 \qquad 1_\gamma\longmapsto\gamma(a),
\]
is a coherent \(\Gamma\)-equivariant equivalence of underlying
\(S\)-modules: it is the induced map adjoint to the class \(a\), and
it is an isomorphism on every homotopy group.  The diagonal unit maps
to \(\sum_\gamma\gamma(a)=1\).  Equivalently, under the
induction--coinduction adjunction its invariant unit is the specified
unit of \(S_k\).  Localizing gives the same statement for
\(R_n\to R_{n,k}\).

Let an exact functor be defined on the underlying \(K(n)\)-local
\(S\)-modules, as are the inclusion followed by the ordinary tests
used for spectral finiteness below.  It preserves this finite sum, its
action, and the diagonal map.  Finite induction equals finite
coinduction, so the homotopy fixed points of the transformed
\(R_{n,k}\) recover the transformed \(R_n\) and its specified unit.  This argument does not require that the functor
commute with an arbitrary homotopy limit.  On coefficient cochains the
trace-one averaging and contraction of
\Cref{h3:lem:framework-semilinear} likewise give exact semilinear descent,
even when \(p\) divides \(|\Gamma|\).  Finally, scalar extension is
faithful: its underlying finite free module of positive rank detects
a zero cofiber.

\subsection{The ordinary residue test}

\begin{lemma}\label{h3:des:flat-functions}
Put \(\widetilde A=W(k)[[x,y]]\), and let \(Q\) be profinite.
The module \(C_{\hthreects}(Q,\widetilde A)\), for the
\((p,x,y)\)-adic topology, is pro-free and flat over
\(\widetilde A\).  For every ordinary \(E_3(k)\)-module \(N_Q\) whose graded
coefficient module is \(C_{\hthreects}(Q,(E_3(k))_*)\), ordinary
smash with \(E_2(k)\) over
the \(p\)-local sphere \(\mathbb S_{(p)}\) has homotopy
\begin{equation}\label{h3:des:ordinary-kunneth}
 (E_2(k))_*E_3(k)\otimes_{(E_3(k))_*}
 C_{\hthreects}(Q,(E_3(k))_*).
\end{equation}
The first tensor factor in \eqref{h3:des:ordinary-kunneth} denotes
absolute ordinary cooperations over \(\mathbb S_{(p)}\).  This applies
in particular to the actual term \(N_Q=\mathcal N_3^q\) for
\(Q=P_3^q\), with its first-factor module structure.
\end{lemma}

\begin{proof}
For each power \(\mathfrak m^a\) of \(\mathfrak m=(p,x,y)\),
coefficient reduction gives
\[
 C_{\hthreects}(Q,\widetilde A)/\mathfrak m^a
 \cong C_{\mathrm{lc}}(Q,\widetilde A/\mathfrak m^a).
\]
Surjectivity follows by lifting on a finite clopen partition.
For the kernel, partition the monomials in \(p,x,y\) among the
finitely many monomial generators of \(\mathfrak m^a\); the resulting
division operations on coefficients are continuous and write a
kernel function as a sum of these generators times continuous
functions.  The right side is the filtered colimit of finite free
modules associated to finite clopen partitions.  It is flat over
the Artinian local ring \(\widetilde A/\mathfrak m^a\).

Choose a basis of the residue module over \(k\), lift the basis
elements, and let \(F\) be the free \(\widetilde A\)-module on those
lifts.  An element of \(\widehat F\) is a family of coefficients for which,
modulo each \(\mathfrak m^a\), only finitely many are nonzero.
Consequently \(\widehat F/\mathfrak m^a\) is the free
\(\widetilde A/\mathfrak m^a\)-module on the chosen basis; continuous
monomial division gives the kernel equality used here.  The map
\(\widehat F\to C_{\hthreects}(Q,\widetilde A)\) is an isomorphism
modulo \(\mathfrak m\).  It is an isomorphism modulo every
\(\mathfrak m^a\): both modules there are flat, and successive
nilpotent-ideal reduction proves injectivity and surjectivity.
Taking the inverse limit proves pro-freeness.  For the ordinary
flatness conclusion, \(\widetilde A=W(k)[[x,y]]\) is a regular complete
Noetherian local ring, and \(\mathfrak m=(p,x,y)\) is generated by the
minimal regular sequence \(p,x,y\).  The module \(F\) is free, hence
flat, and \(F/\mathfrak mF\) is free over \(\widetilde A/\mathfrak m\).
Thus \cite[Proposition~3.13]{BarthelStapleton2016} applies and makes the
ordinary \(\mathfrak m\)-adic completion \(\widehat F\) flat.  The
pro-free characterizations are also described in
\cite[Theorem~A.9 and Proposition~A.13]{HS99}.

Now use the ordinary module tensor identity
\[
 E_2(k)\wedge N_Q
 \simeq
 (E_2(k)\wedge E_3(k))\otimes_{E_3(k)}N_Q.
\]
The module Künneth spectral sequence has no positive Tor terms
by the flatness just proved.  Periodicity reduces graded modules over this coefficient ring to
parity-graded modules over the regular local ring
\(W(k)[[x,y]]\), of global dimension three.  A projective resolution
of length at most three therefore gives bounded convergence of this
Künneth calculation.
This proves \eqref{h3:des:ordinary-kunneth}.
\end{proof}

\begin{lemma}\label{h3:des:residue-cooperations}
Let
\[
 \overline E_2=E_2(k)/(p,u_1).
\]
Apply the ordinary relative tensor \(\overline E_2\otimes_{S_k}(-)\)
to the \(K(3)\)-completed Amitsur object for \(E_3(k)\).  Its homotopy is
even periodic.  After using the test periodicity generator, the
degree-zero coefficient in cosimplicial degree \(q\) is
\[
 \colim_B C_{\hthreects}(P_3^q,B),
\]
where \(B\) runs through the finite connected-marking algebras
over \(k((y))\) of \Cref{h3:sec:coheight-one}.  The cofaces give their
actual \(P_3\)-action, and the test of the unit is the constant
cochain.
\end{lemma}

\begin{proof}
Before imposing the residue ideal, both Morava theories are
Landweber exact over \(BP\): the sequence
\(p,u_1,\ldots,u_{n-1}\) is regular and the height-\(n\) coefficient
is a unit, and these facts are retained by finite unramified scalar
extension.  Landweber exactness gives
the homology base-change formula for all spectra, as recalled in the
introduction and Example~0.1(d) of \cite{HoveyStrickland}.  Applying it
first to \(E_2(k)_*(E_3(k))\), then to \(E_3(k)_*(BP)\) and using the
symmetry of ordinary smash, gives
\begin{equation}\label{h3:des:absolute-landweber-cooperations}
 (E_2(k))_*E_3(k)
 \cong (E_2(k))_*\otimes_{BP_*}BP_*BP
                 \otimes_{BP_*}(E_3(k))_*.
\end{equation}
The two tensor products use the two unit maps of \(BP_*BP\).
The strict \(p\)-typical formal-law groupoid represented by
\((BP_*,BP_*BP)\), including its units and composition, is recalled in
\cite[Section~2, Lemma~2.9 and the following discussion, p.~7]{BhattacharyaEgger2019}.
Naturality of the two multiplicative orientations identifies these
base-changed operations with those of the oriented formal isomorphisms.
This is an absolute ordinary cooperation calculation,
which is precisely the first factor in
\eqref{h3:des:ordinary-kunneth}.

We also record the needed flatness over each coefficient factor.
For a Landweber-exact \(BP_*\)-algebra \(B\),
\cite[Lemma~2.2]{HoveyStrickland} makes
\(B\otimes_{BP_*}BP_*BP\) flat for its other \(BP_*\)-action.
Hopf-algebroid conjugation gives the analogous assertion for
\(BP_*BP\otimes_{BP_*}B\).  Applying these two assertions with
\(B=(E_2(k))_*\) and \(B=(E_3(k))_*\), and then base changing the
remaining unit, proves that
\eqref{h3:des:absolute-landweber-cooperations} is flat over either
coefficient factor.  Hence the sequence \((p,u_1)\) on the height-two
factor acts regularly on these cooperations.
Tensoring over the height-three factor with the flat module in
\Cref{h3:des:flat-functions} preserves that regularity.  The two
residue cofibers therefore create no odd homotopy in the absolute
smash product.

We next pass from this absolute residue calculation to the
relative test in the statement.  Let
\[
 C=\cofib(\mathbb S_{(p)}\longrightarrow S).
\]
The completion map is a mod-\(p\) equivalence, so \(C/p=0\).
Thus multiplication by \(p\) is invertible on \(C\), and the
\(p\)-local spectrum \(C\) is rational.  The spectrum
\(\overline E_2\) is rationally acyclic, already after its
first residue cofiber.  Hence \(\overline E_2\wedge C=0\).
The unit map \(\overline E_2\to\overline E_2\wedge S\) is
an equivalence of spectra.  Its composite with the action map
\(\overline E_2\wedge S\to\overline E_2\) is the identity.
The action map is therefore an equivalence of right
\(S\)-modules.  Consequently, for every
\(S\)-module \(N\), there is a natural equivalence
\begin{equation}\label{h3:des:absolute-relative-residue}
 \overline E_2\wedge N
 \simeq(\overline E_2\wedge S)\otimes_S N
 \simeq\overline E_2\otimes_S N.
\end{equation}

For \(S_k\)-modules, the two scalar actions on the last
object give an action of
\[
 S_k\otimes_S S_k
       \simeq\prod_{\sigma\in\Gal(k/\F_p)}S_k.
\]
The multiplication map \(S_k\otimes_S S_k\to S_k\) is projection
to the identity factor.  Indeed, balancing the two scalar actions is
\[
 (\overline E_2\otimes_S N)
   \otimes_{S_k\otimes_S S_k}S_k
 \simeq \overline E_2\otimes_{S_k}N.
\]
Thus the identity idempotent selects
\(\overline E_2\otimes_{S_k}N\).  On the absolute
residue coefficient algebra already calculated, this is the
equal-embeddings summand of
\(k\otimes_{\F_p}k\simeq\prod_\sigma k\).
It is a retract and therefore remains even.  Applying this
to \(N=C_{\hthreects}(P_3^q,E_3(k))\) gives the required relative
term.  The comparison \eqref{h3:des:absolute-relative-residue}
and the identity idempotent are natural for \(S_k\)-linear
maps.  Every \(P_3\)-cobar coface is \(S_k\)-linear,
so these identifications commute with all cofaces and
with the constant unit.  Regularity was needed only for
the absolute calculation before the residue cut.

In this identity summand, the coefficients are described
by the formal-group-isomorphism Hopf algebroid over \(k\).
Under an isomorphism of formal groups the ideals
\((p,v_1)\) agree on the two sides.  Thus this cut is \(p=x=0\)
on the height-three factor.  The height-two \(v_2\) is a unit,
so \(y\) is a unit as well.  After fixing the test periodicity generator,
the conversion from graded to ungraded coordinates retains the invertible
leading coefficient of the isomorphism
\cite[(2.10)--(2.11) and (2.15)--(2.16), pp.~7--9]{BhattacharyaEgger2019}.
For the oriented isomorphism used here, write
\(aT^{p^2}\) and \(bT^{p^2}\) for the first nonzero terms of the
two \(p\)-series and \(rT\) for the linear term of the isomorphism.
The equation intertwining the \(p\)-series gives
\(ra=br^{p^2}\), or \(r^{p^2-1}=a/b\).  Both \(a\) and \(b\) are units,
and \(p^2-1\) is prime to \(p\), so this is the finite étale equation
identifying the nonzero height-two coefficients.  The remaining equations
are precisely the coefficients of an isomorphism from the specialized
law to the Honda law.

The finite connected-marking stages of
\Cref{h3:thm:coordinate-comparison,h3:prop:integral-frames} represent
these truncated equations.  Their union represents the full
isomorphism algebra; every polynomial in the cooperation
coordinates belongs to a finite such stage.  Thus the ordinary
algebraic tensor product gives this algebraic union, without its
valuation completion.

It remains to check the parameter topology.  Reduction of
continuous power-series functions modulo \((p,x)\) gives
continuous functions into \(k[[y]]\): coefficient lifting proves
surjectivity, and continuous monomial division by \(p,x\) identifies
the kernel with \((p,x)C_{\hthreects}(Q,\widetilde A)\).
There is an equality
\begin{equation}\label{h3:des:compact-laurent}
 C_{\hthreects}(Q,k[[y]])[y^{-1}]
   = C_{\hthreects}(Q,k((y))).
\end{equation}
Indeed, the image of compact \(Q\) in the local field is bounded,
so multiplication by one power of \(y\) puts that image in
\(k[[y]]\).  The reverse inclusion in
\eqref{h3:des:compact-laurent} is immediate.  A finite étale
algebra over \(k((y))\) is a finite product of finite-dimensional
complete vector spaces over it; a basis gives the identical
assertion at each finite marking stage.  Consequently the final
union is a colimit on continuous cochain complexes with a common
stage for every compact family.

Lastly, all the identifications were made in the isomorphism
Hopf algebroid.  Its composition law transports the tautological
marking and is exactly the coface action.  The identity arrow
gives the constant test-unit.  This proves the assertions about
the maps, not just the coefficient modules.
\end{proof}

\subsection{The determinant map and the height-two basis}

Choose once and for all a topological generator of
\(\Z_p^\times/\mu_{p-1}\), and let \(c_{\mathrm{det}}\) be the reduced-determinant
logarithm normalized to send it to \(1\).  On the standard extended
height-three stabilizer set \(c_{\mathrm{det}}(g,\sigma)=c_{\mathrm{det}}(g)\) for \(g\in P_3\).
The reduced norm is invariant under Galois conjugation, so this is a
surjective continuous homomorphism to \(\Z_p\).  Let \(T^1\) be the
Devinatz--Hopkins fixed-point spectrum for its kernel, and let
\(\gamma\) represent the chosen generator in the residual quotient.
Their Proposition~8.1 gives a fiber sequence with middle map
\(1-\gamma\) and first map the unit.  Negating the target of the middle
map gives the convention used here,
\[
 T\longrightarrow T^1\xrightarrow{\gamma-1}T^1
       \xrightarrow{\partial_+}\Sigma T.
\]
The first map stays the same; if \(\partial_-\) denotes the boundary in
the \(1-\gamma\) sequence, the resulting boundary is
\(\partial_+=-\partial_-\).  By
\cite[Theorem~6 and Propositions~8.1--8.2]{DevinatzHopkins}, the boundary
of the unit is an actual permanent degree-one class detected by
\(\pm c_{\mathrm{det}}\).  We determine its sign in our cobar convention locally.

We use the cochain differential \(\partial=d^0-d^1\) in degree zero,
so \((\partial b)(g)=g b-b\).  For the standard field \(F=\F_{p^3}\),
\cite[Theorem~2(i)--(ii)]{DevinatzHopkins} identifies the Morava module
of \(T^1\) with \(C_{\hthreects}(\Z_p,E_{F,*})\) and its Adams
spectral sequence with continuous cohomology.  Realize this module with
the \(T^1\) factor first and the test \(E_F\) factor second.  If
\(i:T^1\to E_F\) is the standard map and \(h\) represents the coset with
\(c_{\mathrm{det}}(h)=t\), its homogeneous evaluation is
\(\operatorname{ev}_t=\mu\circ(h^{-1}i\otimes1)\).
The \(\Z_p\)-valued submodule \(C_{\hthreects}(\Z_p,\Z_p)\) contains
the unit and the lift needed below.  Write \(t\) for the coordinate
with \(c_{\mathrm{det}}(\gamma)=1\).  The actual evaluation formula gives two commuting
actions: the residual action on the \(T^1\) factor is
\[
 (\gamma_L f)(t)=f(t-1),
\]
while the action induced by the other Morava factor is
\((g_R f)(t)=g(f(t+c_{\mathrm{det}}(g)))\).  On the \(\Z_p\)-valued submodule used
here this is simply
\[
 (g_R f)(t)=f(t+c_{\mathrm{det}}(g)).
\]
Indeed, equivariance of \(i\) gives
\(\operatorname{ev}_t\circ(\gamma\otimes1)=\operatorname{ev}_{t-1}\),
using \(\gamma^{-1}h\) for the new representative, while
\(\operatorname{ev}_t\circ(1\otimes g)
 =g\circ\operatorname{ev}_{t+c_{\mathrm{det}}(g)}\), using \(hg\) for that
representative.  These are identities of the actual multiplication
maps, and the \(\Z_p\) values have trivial coefficient action.
Both displayed Morava modules are even,
so completed Morava homology of the fiber sequence gives the exact
coefficient sequence
\[
 0\longrightarrow E_{F,*}\longrightarrow
 C_{\hthreects}(\Z_p,E_{F,*})
 \xrightarrow{D=\gamma_L-1}
 C_{\hthreects}(\Z_p,E_{F,*})\longrightarrow0,
\]
where the first map includes the constant functions.  We fix the sign of every descent edge used below on its coefficient
Postnikov layer.  For normalized cosimplicial degree \(s\) and internal
cochain degree \(r=-t\), use total degree \(s+r\) and differential
\[
 d_{\mathrm{Tot}}=d_{\mathrm{cos}}+(-1)^s d_{\mathrm{int}},
 \qquad d_{\mathrm{cos}}=\sum_i(-1)^i d^i.
\]
Thus the unit edge is the constant \(1\), and the internal-degree-zero
cochain differential is the one just specified.  Model suspension by
\(K[1]^n=K^{n+1}\) with differential \(-d_K\), with the strict inverse
shift for desuspension.  The comparison
\(\operatorname{Tot}(K[1])\to(\operatorname{Tot}K)[1]\) is
multiplication by \((-1)^s\) in column \(s\).

Here is the resulting boundary sign.  For a degreewise exact sequence
of normalized cochain complexes \(0\to A\xrightarrow{i}B\xrightarrow{D}C\to0\),
use
\[
 \operatorname{Cone}(i)^n=B^n\oplus A^{n+1},\qquad
 d(b,a)=(d_Bb+i(a),-d_Aa),
\]
with inclusion \(j(b)=(b,0)\) and projection \(p(b,a)=a\) to \(A[1]\).
This cone has the rotation convention of Section~2: the map
\(a\mapsto((0,a),-i(a))\) from \(A[1]\) to \(\operatorname{Cone}(j)\)
identifies the next cone projection with \(-i[1]\).
The quotient \(q:\operatorname{Cone}(i)\to C\), \(q(b,a)=D(b)\), is a
quasi-isomorphism under \(B\), so the boundary is the roof \(pq^{-1}\).
For a cocycle \(c_0=D(b)\), put \(a_1=i^{-1}d_Bb\).  The cone cocycle
lifting \(c_0\) is \((b,-a_1)\), and its boundary is \(-a_1\).
In particular this boundary is the negative of the short-exact-cochain
connecting cocycle \(+a_1\).

This is also the sign of the actual desuspended boundary at its first
Adams edge.  On the coefficient Postnikov layer, the comparison from
the totalization of a termwise cone to the cone of the totalized map is
\((b,a)\mapsto(b,(-1)^s a)\) in column \(s\).  It commutes with the cone
projection and the displayed suspension comparison.  Evenness gives
the displayed short exact coefficient sequence, and the unit's boundary
first occurs in bidegree \((1,0)\); hence this layer computes its leading
edge.  The strict inverse shift contributes no further desuspension sign.

For the constant unit, take \(b(t)=-t\).  Then
\[
 (Db)(t)=b(t-1)-b(t)=1,\qquad
 (d_{\mathrm{cos}}b)(g)(t)=b(t+c_{\mathrm{det}}(g))-b(t)=-c_{\mathrm{det}}(g).
\]
The boundary-unit map \(\Sigma^{-1}(\partial_+\circ1)\) is therefore
detected by \(-d_{\mathrm{cos}}b=+c_{\mathrm{det}}\).  Define
\[
 \zeta=\Sigma^{-1}(\partial_+\circ1)\in\pi_{-1}T.
\]
It is an actual sphere map detected by the inflated normalized logarithm
\(c_{\mathrm{det}}\).  The normalization also agrees with the two-term group-ring
resolution: its degree-one generator maps to the bar generator
\([\gamma]\), whose boundary is \(\gamma-1\), so the corresponding
one-cocycle takes value \(1\) on \(\gamma\).
This normalization is used for all later descent edges.  On the
internal-degree-zero rows the total-complex Alexander--Whitney product
has no interchange sign, so these edges retain the usual cup products
and their naturality under the Moore quotient maps.
The relative coface
calculation restricts this class to the same logarithm on \(P_3\), and
finite-field descent retains this map and the unit.

\begin{theorem}\label{h3:thm:height-two-test}
Set
\[
 W=L_2S\vee\Sigma^{-1}L_2S,\qquad
 w=(1,\zeta):W\longrightarrow X.
\]
The first component is the canonical unit, and \(L_{K(2)}w\)
is an equivalence.  More explicitly, the specified maps give
\[
 (1,\zeta):R_2\vee\Sigma^{-1}R_2
       \xrightarrow{\ \simeq\ }L_{K(2)}T.
\]
\end{theorem}

\begin{proof}
The unit and the integral class just constructed are maps out
of \(S\).  Thus they define \(w\) by applying \(L_2\); no
extension of a map defined only on \(R_2\) is being asserted.

Use \Cref{h3:prop:exact-descent,h3:des:residue-cooperations} over
\(S_k\).  By \Cref{h3:thm:coheight-one-comparison}, the actual
coefficient augmentation gives
\[
 R\Gamma_{\hthreects}(H',L)\simeq k.
\]
It retains compact parameters and the action of \(P_3/H'\).
This quotient is \(\Z_p\), with trivial action on \(k\).
The continuous group-ring resolution
\[
 0\longrightarrow k[[\Z_p]]
   \xrightarrow{\gamma-1}k[[\Z_p]]
   \longrightarrow k\longrightarrow0
\]
therefore shows that the residue-tested spectral sequence has
only columns zero and one.  They are generated by the constant
unit and the normalized logarithm.  There is no possible
differential, and the specified maps \(1,\zeta\) detect these
two generators by the coface identification in
\Cref{h3:des:residue-cooperations}.

Because \(L_2\) is smashing and \(\overline E_2\) is
\(E(2)\)-local, the ordinary relative tests satisfy
\[
 \overline E_2\otimes_{S_k}(W\otimes_S S_k)
       \simeq \overline E_2\vee\Sigma^{-1}\overline E_2,
 \qquad
 \overline E_2\otimes_{S_k}(X\otimes_S S_k)
       \simeq \overline E_2\otimes_{S_k}R_{3,k}.
\]
The preceding two detected classes therefore prove that the test of
\(w\otimes_S S_k\) is an equivalence.  By associativity this test is
\(\overline E_2\otimes_S w\), and
\eqref{h3:des:absolute-relative-residue} identifies it with the
absolute ordinary residue test of \(w\).
The residue theory \(\overline E_2\) has Bousfield class \(K(2)\);
its coefficient ring is a graded field after a finite field
extension.  Consequently \(w\) is a \(K(2)\)-equivalence.
Applying \(L_{K(2)}\) to \(w\) therefore gives an equivalence with
source \(R_2\vee\Sigma^{-1}R_2\) and target \(L_{K(2)}T\).
The scalar passage preserved the two original \(S\)-maps by the
naturality of the test and the unit-compatible finite-field descent.
Its components are consequently the specified unit and \(\zeta\).
\end{proof}
